\begin{table}[htbp]\centering\small
\caption{Style-adjusted comparisons per ten matches}\label{tab:headlines}
\resizebox{\textwidth}{!}{\begin{tabular}{llrr}\toprule
Target & Outcome & Gap [95 percent CI] & Holm p \\
\midrule
Barcelona & Fewer yellow cards received & 1.29 [-0.96, 3.55] & 1.000 \\
Barcelona & More opponent yellow cards & -0.75 [-2.37, 0.88] & 1.000 \\
Barcelona & Net yellow cards & 0.55 [-1.85, 2.94] & 1.000 \\
Barcelona & Net fouls & 3.26 [-7.98, 14.50] & 1.000 \\
Barcelona & Net red cards & -0.36 [-1.12, 0.41] & 1.000 \\
Barcelona & Net penalties & -0.10 [-0.95, 0.75] & 1.000 \\
Real Madrid & Net yellow cards & -1.95 [-4.35, 0.45] & 0.941 \\
Real Madrid & Net fouls & 12.89 [2.25, 23.54] & 0.203 \\
Real Madrid & Net penalties & -0.64 [-1.61, 0.33] & 1.000 \\
Real Madrid & Net red cards & 0.40 [-0.42, 1.21] & 1.000 \\
Real Madrid & Net yellow cards, given called fouls & -3.07 [-4.82, -1.33] & 0.018 \\
\bottomrule\end{tabular}}
\par\medskip\begin{minipage}{\textwidth}\footnotesize
Barcelona is compared with Real Madrid in La Liga; Real Madrid is compared with other elite clubs excluding Barcelona in the Champions League. Positive values favor the target, including the sign-reversed own-yellow outcome. Intervals are pointwise season-block clustered intervals. Holm p-values adjust the family of eleven displayed endpoints. These are conditional associations, not estimates of incorrect decisions.
\end{minipage}\end{table}
